\documentclass[11pt]{article}
\usepackage[margin=1in]{geometry}
\usepackage{amsmath,amssymb,amsthm,mathtools}
\usepackage{graphicx}
\usepackage{xcolor}
\usepackage{microtype}
\usepackage[colorlinks=true,linkcolor=teal!70!black,citecolor=teal!70!black,urlcolor=teal!70!black]{hyperref}

\newtheorem{theorem}{Theorem}
\newtheorem{corollary}[theorem]{Corollary}
\newcommand{\R}{\mathbb R}
\newcommand{\Besov}[1]{B^{#1}_{\infty,\infty}}
\newcommand{\Diff}{\Delta}

\title{Power Majorants in the Zygmund Whitney Problem:\\
H\"older Cases and Collapse to One Endpoint}
\author{Run \texttt{fa\_banach\_001}, \texttt{agent\_lane\_14}}
\date{11 August 2026}

\begin{document}
\maketitle

\section*{Status and source question}

\textbf{Status:} candidate substantial structural partial result, likely
valid.  Shvartsman~\cite{Shvartsman2009} defines
$C^k\Lambda^m_\omega(\R^n)$ by bounded derivatives through order $k$ and
the estimate
\[
 \|\Diff_h^mD^\alpha f\|_\infty\leq C\omega(|h|),
 \qquad |\alpha|=k.
\]
On source PDF page~3, Problem~1.2 asks whether these spaces possess the
finiteness property when $m>2$, or when $m=2$ and $k>0$.

\begin{center}
\includegraphics[width=\textwidth]{figures/open_problem_crop.png}
\end{center}

The result below treats all power majorants $\omega(t)=t^s$ with $0<s<m$.
It settles every noninteger power and shows that every integer power is not a
new higher-order problem: it is exactly the classical $m=2$, $\omega(t)=t$
endpoint with a shifted derivative index.

\section*{Proof intuition}

For smoothness $s$ strictly below the chosen difference order $m$, the
quantity $\sup_h |h|^{-s}\|\Diff_h^m g\|_\infty$ is precisely the usual
finite-difference norm of the H\"older--Zygmund/Besov space
$\Besov{s}$.  Taking $k$ bounded classical derivatives simply shifts Besov
smoothness from $s$ to $k+s$.  If $s$ is nonintegral this is an ordinary
H\"older space, where the finiteness theorem is known.  If $s=r$ is integral,
the same space can instead be written using $k+r-1$ derivatives followed by
one Zygmund derivative; this is $C^{k+r-1}\Lambda_t^2$.  The unresolved part
is therefore an endpoint selection problem, not a proliferation of distinct
higher-difference problems.

\section*{Structural partial theorem}

We use inhomogeneous bounded-function norms throughout.

\begin{theorem}[Power-scale classification]\label{thm:classification}
Let $n\geq1$, $k\geq0$, $m\geq2$, and $0<s<m$.  Put
$\omega_s(t)=t^s$.  Then
\begin{equation}\label{eq:main-equivalence}
 C^k\Lambda^m_{\omega_s}(\R^n)=\Besov{k+s}(\R^n)
\end{equation}
with equivalent norms, with constants depending only on $n,k,m,s$.

If $s=a+\theta$ with $a\in\{0,\ldots,m-1\}$ and $0<\theta<1$, then
\begin{equation}\label{eq:holder}
 C^k\Lambda^m_{t^s}(\R^n)=C^{k+a,\theta}(\R^n),
\end{equation}
again with equivalent norms.  Consequently this space has the finiteness
property.

If $s=r\in\{1,\ldots,m-1\}$ is integral, then
\begin{equation}\label{eq:endpoint}
 C^k\Lambda^m_{t^r}(\R^n)
 =C^{k+r-1}\Lambda^2_t(\R^n)
\end{equation}
with equivalent norms.  Hence every integer power-majorant instance of
Problem~1.2 is equivalent to the classical Zygmund endpoint family.
\end{theorem}

\begin{proof}
The standard finite-difference characterization of inhomogeneous Besov
spaces~\cite[finite-difference and lifting theorems]{Triebel1983} states that,
whenever $0<s<m$,
\begin{equation}\label{eq:difference}
 \|g\|_{\Besov{s}}
 \asymp \|g\|_\infty+
 \sup_{0<|h|\leq1}\frac{\|\Diff_h^m g\|_\infty}{|h|^s}.
\end{equation}
The restriction $|h|\leq1$ is immaterial for the source norm, since
$\|\Diff_h^m g\|_\infty\leq2^m\|g\|_\infty$ for every $h$.

The same theory gives the derivative-lifting equivalence
\begin{equation}\label{eq:lifting}
 \|f\|_{\Besov{k+s}}
 \asymp
 \sum_{|\beta|\leq k}\|D^\beta f\|_\infty
 +\sum_{|\alpha|=k}\|D^\alpha f\|_{\Besov{s}}
\end{equation}
on $C_b^k(\R^n)$.  Substitution of \eqref{eq:difference} into
\eqref{eq:lifting} is exactly the norm in the definition of
$C^k\Lambda^m_{t^s}$, proving \eqref{eq:main-equivalence}.

For nonintegral $s=a+\theta$, the classical identification
$\Besov{k+a+\theta}=C^{k+a,\theta}$ proves \eqref{eq:holder}.  The finiteness
property for bounded $C^{q,\theta}$ spaces is part of the Whitney finiteness
theory quoted in the source introduction and developed in
Fefferman's work~\cite{Fefferman2005,Fefferman2006}.  Equivalent norms
transfer both the local hypothesis and the global conclusion, so the
finiteness property transfers to $C^k\Lambda^m_{t^s}$.

Finally let $s=r$ be integral and put $q=k+r-1$.  Applying
\eqref{eq:main-equivalence} first to $(k,m,s)=(k,m,r)$ and then to
$(k,m,s)=(q,2,1)$ gives
\[
 C^k\Lambda^m_{t^r}=\Besov{k+r}
 =\Besov{q+1}=C^q\Lambda^2_t,
\]
which is \eqref{eq:endpoint}.
\end{proof}

\begin{corollary}[The named $Z_m$ scale]
For every $m\geq2$,
\[
 Z_m(\R^n):=\Lambda^m_{t^{m-1}}(\R^n)
 =C^{m-2}\Lambda^2_t(\R^n)
\]
with equivalent norms.  Thus a positive answer for the source's $m=2$
endpoint for every derivative index would settle all of its named spaces
$Z_m$.
\end{corollary}

\section*{Verification, limitations, and upgrade obstruction}

The proof uses the finite-difference characterization only in its valid
strict range $m>s$.  At an integer $s$ it does \emph{not} identify the space
with a Lipschitz class; instead it compares two difference descriptions of
the same endpoint Besov space, which is why \eqref{eq:endpoint} is safe.

This does not settle arbitrary majorants, the borderline case $s=m$, or the
integer endpoint $C^q\Lambda_t^2$ for $q>0$.  A direct check of the later sharp
Lipschitz-selection theorem found that its range must lie in a Banach space,
whereas the source's range is the nonlinear hyperbolic jet space
$(\mathcal P_L\times\mathcal K,d_\omega)$.  A Kuratowski embedding does not
preserve convexity of the polynomial fibers.  A second endpoint route would
be a uniformly bounded-cardinality molecular decomposition of finite sums of
point evaluations in a $B^{-q-1}_{1,1}$ predual; that decomposition remains
unproved here.

Exact-title, arXiv-id, exact-problem-phrase, generalized-Zygmund, and
hyperbolic-jet-space searches through 11 August 2026 found no later full
resolution.  The equivalences used above are classical, so novelty confidence
for the structural application is low.

\section*{Human-review recommendation}

Verify the inhomogeneous bounded-function convention in
\eqref{eq:lifting}, and confirm that the selected $C^{q,\theta}$ finiteness
formulation has constants independent of the closed trace set.  After that,
the parameter shift in \eqref{eq:endpoint} is formal.  For a full upgrade,
focus on the nonlinear-bundle finiteness lemma or the endpoint molecular
decomposition, not on additional difference-order manipulations.

\begin{thebibliography}{9}
\bibitem{Shvartsman2009}
P.~Shvartsman,
\emph{The Whitney extension problem for Zygmund spaces and Lipschitz
selections in hyperbolic jet-spaces}, arXiv:0905.2602 (2009).

\bibitem{Triebel1983}
H.~Triebel,
\emph{Theory of Function Spaces}, Monographs in Mathematics 78,
Birkh\"auser, 1983.

\bibitem{Fefferman2005}
C.~Fefferman,
\emph{A sharp form of Whitney's extension theorem},
Ann. of Math. (2) 161 (2005), 509--577.

\bibitem{Fefferman2006}
C.~Fefferman,
\emph{Whitney's extension problem for $C^m$},
Ann. of Math. (2) 164 (2006), 313--359.
\end{thebibliography}

\end{document}

